% ============================================================
% currency.tex — one framework for money and numbers.
% Language-aware: follows \InvoiceLanguage unless \CurrencyLocale says otherwise.
%
%   \money{EUR}{1200.50}        EUR 1,200.50      (de: 1.200,50 EUR / EUR 1.200,50)
%   \money{USD}{3.2}[\mn]       USD 3.2m          (de: USD 3,2 Mio.)
%   \money{USDC}{250000}        USDC 250,000
%   \money{EUR}{0.004}[/kWh]    EUR 0.004/kWh     (suffix printed verbatim)
%   \pct{12.5}                  12.5 %            \num{1200}  1,200
%   \DeclareCurrency{GBP}{\pounds}{fiat}          add a currency
%
% Type amounts with a DOT decimal and no grouping (1200.50); the locale
% decides how it renders. Choices come from style/defaults.sty:
%   \CurrencyStyle    iso-before | symbol | iso-after
%   \CurrencyLocale   en | de
%   \CurrencyCodeFont font switch for the code (e.g. {} or \sffamily)
% Tables: put the code in the header ("Price (EUR)") and use siunitx
% S columns for the cells: S[table-format=4.2].
% ============================================================
\makeatletter
\RequirePackage{siunitx}
\RequirePackage{xparse}
\RequirePackage{etoolbox}
\RequirePackage[official]{eurosym}
\providecommand{\CurrencyStyle}{iso-before}
\providecommand{\CurrencyLocale}{en}
\providecommand{\CurrencyCodeFont}{}

% ---------- Registry: code -> symbol, kind (fiat | token) ----------
\newcommand{\DeclareCurrency}[3]{%
  \csdef{cur@sym@#1}{#2}\csdef{cur@kind@#1}{#3}}
\DeclareCurrency{EUR}{\officialeuro}{fiat}
\DeclareCurrency{USD}{\$}{fiat}
\DeclareCurrency{GBP}{\pounds}{fiat}
\DeclareCurrency{CHF}{}{fiat}
\DeclareCurrency{JPY}{\textyen}{fiat}
\DeclareCurrency{USDC}{}{token}
\DeclareCurrency{USDT}{}{token}
\DeclareCurrency{EURC}{}{token}
\DeclareCurrency{BTC}{}{token}
\DeclareCurrency{ETH}{}{token}

\newcommand{\CurrencyCode}[1]{{\CurrencyCodeFont #1}}
\newcommand{\CurrencySymbol}[1]{\csuse{cur@sym@#1}}
\newcommand{\CurrencyKind}[1]{\csuse{cur@kind@#1}}

% ---------- Number formatting per locale ----------
\newcommand{\cur@setlocale}{%
  \ifdefstring{\CurrencyLocale}{de}{%
    \sisetup{output-decimal-marker={,}, group-separator={.}, group-minimum-digits=4, group-digits=integer}}{%
    \sisetup{output-decimal-marker={.}, group-separator={,}, group-minimum-digits=4, group-digits=integer}}}
\sisetup{input-decimal-markers={.}, mode=match, per-mode=symbol,
  group-minimum-digits=4, group-digits=integer}
\AtBeginDocument{\cur@setlocale}

% Localised magnitudes (use inside the optional suffix)
\newcommand{\mn}{\ifdefstring{\CurrencyLocale}{de}{~Mio.}{m}}
\newcommand{\bn}{\ifdefstring{\CurrencyLocale}{de}{~Mrd.}{bn}}
\newcommand{\kk}{\ifdefstring{\CurrencyLocale}{de}{~Tsd.}{k}}

% ---------- \money{CODE}{amount}[suffix] ----------
\NewDocumentCommand{\money}{m m o}{%
  \begingroup\cur@setlocale
  \ifmmode\text{\cur@render{#1}{#2}{\IfNoValueTF{#3}{}{#3}}}%
  \else\cur@render{#1}{#2}{\IfNoValueTF{#3}{}{#3}}\fi
  \endgroup}
\newcommand{\cur@render}[3]{%
  \ifdefstring{\CurrencyStyle}{iso-after}%
    {\num{#2}#3\nobreakspace\CurrencyCode{#1}}%
    {\ifdefstring{\CurrencyStyle}{symbol}%
      {\ifcsempty{cur@sym@#1}%
        {\CurrencyCode{#1}\nobreakspace\num{#2}#3}%
        {\CurrencySymbol{#1}\num{#2}#3}}%
      {\CurrencyCode{#1}\nobreakspace\num{#2}#3}}}

% ---------- Percentages ----------
\NewDocumentCommand{\pct}{m}{\begingroup\cur@setlocale
  \ifdefstring{\CurrencyLocale}{de}{\num{#1}\,\%}{\num{#1}\%}\endgroup}
\makeatother
